%
% Portada.
%

% Nota: Sería más cómodo emplear el comando \maketitle que genera una portada de forma automática, pero
% no incluye toda la información que es necesario incluir en la memoria de un proyecto de fin de carrera
% de la Facultad de Informática de A Coruña.
%

\begin{titlepage}

  \begin{center}

    % Logotipo de la universidad.
    \includegraphics[width=6cm]{./eps/logo_udc.eps}
    \vspace{0.4cm}

    % Nombre de la facultad, de la universidad y del departamento en que se realiza el PFC.
    {\Large{\textbf{Universidade da Coruña}}} \\
    {\it \large{\textbf{Mestrado en Profesorado de Educación Secundaria Obrigatoria e Bacharelato, Formación Profesional e Ensino de Idiomas}}}
    \vspace{2cm}

    {\Huge Unidade Didáctica: ``Figuras xeométricas planas. Concepto e propiedades''} \\\vspace{0.4cm}
    {\large Unidad Didáctica: ``Figuras geométricas planas. Concepto y propiedades''} \\ \vspace{0.15cm}
    {\large Didactic Unit: ``Plane geometric figures. Concept and properties''}


    % Indicamos el nombre de la titulación oficial que hemos cursado con tanto esfuerzo.
  \end{center}

  \begin{bottompar}
    \begin{flushright}
      \begin{tabular}{rl}
        \large{\textbf{Alumno:}}
        & \large{Chavarría Teijeiro, Marcos} \\

        & \small{(DNI: 33558386-Y)} \\

        \large{\textbf{Especialidade:}}
        & \large{Matemáticas} \\

        & \small{(Itinerario: Tecnoloxía)} \\


        \large{\textbf{Titora (Universidade):}}
        & \large{Pérez González, Mercedes} \\

        \large{\textbf{Centro:}}
        & \large{IES. Leiras Pulpeiro} \\

        & \small{(Enderezo: Rúa da Orquídea, 45, Lugo)} \\

        \large{\textbf{Data:}}
        & \large{14 de Xuño de 2016} \\
      \end{tabular}
    \end{flushright}
  \end{bottompar}

\end{titlepage}
